% RnE 보고서 양식 스타일 파일
% 패키지·환경·명령어 설정만 — 본문 내용 없음

%-- 용지 설정
\usepackage{fapapersize}
\usefastocksize{210mm,297mm}
\usefapapersize{210mm,297mm,25mm,25mm,25mm,25mm}

%-- 수식
\usepackage{mathtools}
\usepackage{amsthm}
\usepackage{amssymb}
\usepackage{mathrsfs}

%-- 그래픽
\usepackage{graphicx}
\usepackage{xcolor}

%-- 텍스트·레이아웃
\usepackage{setspace}
\usepackage{lipsum}
\usepackage{jiwonlipsum}
\usepackage{ragged2e}
\usepackage{textpos}
\usepackage{kskorsymb}

%-- 표
\usepackage{tabularray}
\UseTblrLibrary{booktabs}
\newlength\mycellwidth  % 셀 병합에서 사용하는 길이 변수

%-- 캡션
\usepackage{caption}
\DeclareCaptionLabelFormat{cjk-single-bracket}{[#1 #2]}
\captionsetup[figure]{labelsep=space, labelfont=bf,
  labelformat=cjk-single-bracket, name={Fig.}}
\captionsetup[table]{labelsep=period, labelfont=bf, name={Table}}

%-- 섹션 번호 형식
\renewcommand\thesection{\Roman{section}}
\renewcommand\thesubsection{\arabic{subsection}}
\renewcommand\thesubsubsection{\thesubsection.\arabic{subsubsection}}
\setsecnumformat{\csname the#1\endcsname.\ }
\setcounter{secnumdepth}{3}
\renewcommand{\abstractname}{Abstract}

%-- 정리류 환경
\newtheorem{theorem}{\bfseries\sffamily 정리}
\newtheorem{lemma}[theorem]{도움정리}
\newtheorem{corollary}[theorem]{따름정리}
\theoremstyle{definition}
\newtheorem{definition}[theorem]{정의}
\newtheorem{example}[theorem]{보기}
\theoremstyle{remark}
\newtheorem{remark}[theorem]{참고}
\numberwithin{equation}{section}
\renewcommand*{\proofname}{증명}

\allowdisplaybreaks

%-- 커스텀 명령어
% 수식 안 쉼표에서 줄바꿈 허용
\newcommand{\splitatcommas}[1]{%
  \begingroup
  \begingroup\lccode`~=`, \lowercase{\endgroup
    \edef~{\mathchar\the\mathcode`, \penalty0 \noexpand\hspace{0pt plus 1em}}%
  }\mathcode`,="8000 #1%
  \endgroup
}

% 체크된 체크박스
\newcommand{\mycheckedbox}{$\square$\!\!\!\!\checkmark}

% 각 페이지 상단 헤더 (1·2·3페이지 섹션 제목)
% 선택 인수: 헤더 아래 수직 여백 (기본 23pt)
\newcommand{\RnEpageheader}[2][23pt]{%
  \begin{textblock*}{\textwidth}(0cm,0cm)%
    \noindent\LARGE\bfseries #2%
  \end{textblock*}%
  \vspace{#1}%
}

%-- 참고문헌 설정
% 참고문헌 설정 (biblatex + biber)

\usepackage[backend=biber, natbib=true, sorting=none,
  citestyle=numeric-comp, bibstyle=numeric,
  maxbibnames=10, minbibnames=6,
  maxcitenames=6, mincitenames=1]{biblatex}
\addbibresource{refs.bib}

\renewbibmacro{in:}{}
\DeclareFieldFormat[article]{volume}{\emph{#1}}
\DeclareFieldFormat{pages}{#1}
\renewcommand{\finalnamedelim}{\ \&\ }

\renewbibmacro*{volume+number+eid}{%
  \printfield{volume}%
  \printfield{number}%
  \setunit{\addcomma\space}%
  \printfield{eid}}
\DeclareFieldFormat[article]{number}{\mkbibparens{#1}}

\AtEveryBibitem{%
  \DeclareFieldFormat[article,inproceedings]{title}{#1}%
  \ifkeyword{kobib}{%
    \renewcommand{\multinamedelim}{, }%
    \renewcommand{\finalnamedelim}{, }%
    \DeclareFieldFormat[article]{volume}{\textbf{#1}}%
    \DeclareFieldFormat*{journaltitle}{#1,}%
    \DeclareFieldFormat*{titlecase}{%
      \ifdef{\currentfield}
        {\ifcurrentfield{journaltitle}
           {\usefield{\textbf}{\currentfield}}%
           {#1}}
        {#1}}%
    \DeclareFieldFormat*[book,thesis]{title}{#1}%
    \DeclareFieldFormat*[book,thesis]{titlecase}{%
      \ifdef{\currentfield}
        {\ifcurrentfield{title}
           {\usefield{\textbf}{\currentfield}}%
           {#1}}
        {#1}}%
  }{}%
}

\AtEveryCitekey{%
  \ifkeyword{kobib}{%
    \renewcommand{\multinamedelim}{, }%
    \renewcommand{\finalnamedelim}{, }%
  }{}%
}

\DefineBibliographyStrings{english}{phdthesis = {박사학위논문}}
\DefineBibliographyStrings{english}{mathesis  = {석사학위논문}}

\newbibmacro*{institution+thesistype+location+date}{%
  \printlist{location}%
  \iflistundef{institution}
    {\setunit*{\addcomma\space}}
    {\setunit*{\addcolon\space}}%
  \printlist{institution}%
  \setunit*{\addspace}%
  \printfield{type}%
  \setunit*{\addcomma\space}%
  \usebibmacro{date}%
  \newunit}

\usepackage{xpatch}
\xpatchbibdriver{thesis}
  {\printfield{type}%
   \newunit
   \usebibmacro{institution+location+date}}
  {\usebibmacro{institution+thesistype+location+date}}
  {}{}

\DeclareDelimFormat[textcite]{nameyeardelim}{}


%-- 요약문 매크로 (학생 편집: 03_summaries.tex)
%============================================================
% 03_summaries.tex — 1·2페이지 표 내용
%
% 이 파일의 각 \newcommand가 다음 칸을 채웁니다.
%   1페이지(연구결과 요약문): summaryGoal / summaryResult / summaryEffect
%   2페이지(과제 수행 과정): summaryExpert / summaryProblemSolving /
%                            summaryFirstMember / summarySecondMember / summaryLearning
%
% \jiwon[숫자]는 한글 더미 텍스트입니다 (Lorem ipsum의 한글 버전).
% 각 칸의 \jiwon[...] 줄을 지우고 실제 내용을 작성하세요.
% 여러 문단을 써도 되고, 줄바꿈도 사용할 수 있습니다.
%============================================================

%%% 1페이지 ─────────────────────────────────────────────────

\newcommand{\summaryGoal}{%
  % '연구 필요성 및 목적' 칸
  \jiwon[1] % ← 이 줄을 지우고 내용을 작성하세요
}

\newcommand{\summaryResult}{%
  % '연구내용 및 연구결과' 칸
  \jiwon[2-3] % ← 이 줄을 지우고 내용을 작성하세요
}

\newcommand{\summaryEffect}{%
  % '연구 기대효과' 칸
  \jiwon[4] % ← 이 줄을 지우고 내용을 작성하세요
}

%%% 2페이지 ─────────────────────────────────────────────────

\newcommand{\summaryExpert}{%
  % '전문가 피드백 활용 과정' 칸
  % 전문가 피드백을 받지 않았다면 아래 한 줄로 대체하세요.
  %   해 당 없 음
  \jiwon[5] % ← 이 줄을 지우고 내용을 작성하세요
}

\newcommand{\summaryProblemSolving}{%
  % '문제해결 과정' 칸 — 팀 전체의 과정을 서술
  \jiwon[6-7] % ← 이 줄을 지우고 내용을 작성하세요
}

\newcommand{\summaryFirstMember}{%
  % '팀원1' 칸 — 00_info.tex의 \myauthorfirstko(저자1)의 역할
  \jiwon[8] % ← 이 줄을 지우고 내용을 작성하세요
}

\newcommand{\summarySecondMember}{%
  % '팀원2' 칸 — 00_info.tex의 \myauthorsecondko(저자2)의 역할
  \jiwon[8] % ← 이 줄을 지우고 내용을 작성하세요
}

\newcommand{\summaryLearning}{%
  % '과제 수행 과정에서의 학습' 칸
  \jiwon[9] % ← 이 줄을 지우고 내용을 작성하세요
}

